Notre Seigneur et Rédempteur parle dans son Évangile
		tantôt avec des mots, tantôt par des faits;
	tantôt il dit une chose avec des mots et une autre par des faits,
	tantôt il dit la même chose par les mots et par les faits.
Ainsi, frères, dans l'Évangile, vous entendez deux histoires:
		celle du figuier stérile et celle de la femme courbée,
	et dans l'une et l'autre est prodiguée de la bonté.
Là, il a parlé en parabole, ici, il s'est manifesté en agissant.
Mais le figuier stérile a la même signification que la femme courbée,
	et le figuier épargné que la femme redressée.
